\section{What changes and what survives}

The correction changes the interpretation of the threshold results, not the underlying equilibrium construction. First, the individual trigger associated with Observation 1 cannot be treated as $\gamma$-invariant over asymmetric spillover configurations. The correct object is the zero set $D_i=0$, whose location can vary with $\gamma$. The symmetric formula in Eq. \eqref{eq:symcross} remains useful, but as a special crossing rather than a global threshold characterization. The reproduced Table 1 benchmarks need not be discarded, but they cannot carry the general invariance interpretation assigned to them.

Second, the aggregate boundary associated with Observation 4 cannot generally be represented by a scalar threshold in $\beta_1+\beta_2$. At $\lambda=0$, the straight line $\beta_1+\beta_2=1$ intersects the corrected equality set at the symmetric point but fails at regular asymmetric points. Figure \ref{fig:aggregate} visualizes this distinction. More generally, Proposition \ref{prop:aggregate} shows that the corrected boundary can vary with $\gamma$ as well as with spillover asymmetry. Accordingly, the published straight boundaries should not be used as global representations of the asymmetric equality set; the relevant object is $D_A=0$ itself.

Third, the fixed average-spillover classification in Observation 5 and the corresponding cutoff equations require qualification. Proposition \ref{prop:misclass} supplies an exact admissible parameter point for which the published cutoff and the actual aggregate R\&D ranking disagree, and continuity makes the disagreement positive-measure. This is a correction of a parameter-region classification. It should not be restated as a reversal of social welfare: the mathematical object corrected here is the aggregate R\&D comparison. Nothing in Propositions 1--3 derives a separate welfare function or establishes that a welfare ranking changes at the counterexample.

Several parts of Ishikawa and Shibata (2021) are unaffected by this correction. The output-stage equilibrium and the closed-form noncooperative and cooperative R\&D solutions are retained. The displayed $\gamma=50$ trigger benchmarks in Table 1 reproduce numerically. The correction also does not overturn the paper's original Propositions 1 and 2 under the dependency considered here. Original Proposition 2 is compatible with the more general point that the relevant comparison may depend on $\beta_1,\beta_2,\gamma$, and $\lambda$; what fails later is the reduction of that dependence to fixed scalar cutoffs. Recomputing the branch choice underlying Eq. (47) shifts the classification locally near the corrected aggregate boundary, but the qualitative pattern summarized in Observation 6 survives at the published calibration.

Finally, the paper's broader statement that stronger market competitiveness tends to enlarge the region in which R\&D competition is preferred is not shown here to reverse globally. What changes is the fixed-cutoff rationale used to summarize that relationship over asymmetric spillovers. Readers who use the published cutoffs to classify asymmetric parameter regions can obtain the wrong regime, whereas readers who use the closed-form equilibria directly retain valid inputs. The practical repair is therefore: retain the equilibria, replace the threshold map by the relevant zero set, and preserve only those qualitative statements that do not require the scalar reduction.

\paragraph{Reproducibility.}
All algebraic identities, exact rational evaluations, and regularity certificates used in Propositions \ref{prop:individual}--\ref{prop:misclass} are verified by deterministic Python/SymPy code. No external empirical dataset and no stochastic simulation are required for any proposition. A reviewer-safe reproducibility package containing the verification scripts, dependencies, and expected outputs accompanies the submission package. Generative-AI assistance was used in developing and reviewing parts of the symbolic-verification and publication workflow; the reported mathematical identities and certificates are encoded as deterministic assertions and remain the responsibility of the author.
